%8. Explaining the lack of prominence: 
%===============================================================================================
% - not a lack of sincerity (petition, list experiment, conjoint analyses)
%===============================================================================================

%LB: ici le paragraphe introductif manquant en reprenant tes commentaires de la version passée. 
The evidence reviewed above gives rise to a preference--prominence puzzle: globally redistributive policies attract majority acceptance when surveys present them to citizens, yet they remain largely absent from public debate and political programmes. \citet{fabre_majority_2025} and \citet{fabre_public_2026} examine several explanations for this discrepancy. Petition and list-experiment results assess whether insincerity or social-desirability bias inflates responses to direct survey questions. Conjoint analyses test whether support persists when respondents choose among competing policy bundles or candidates. Evidence from the Green New Deal illustrates how an adversarial political campaign can erode support. Measures of second-order beliefs identify pluralistic ignorance, while responses to unprompted open-ended questions reveal the low salience of global redistribution. Finally, an election-framing experiment tests whether political institutions organized at the national level direct citizens' attention toward national constituencies and policy responses. Taken together, these results provide little evidence of widespread preference falsification. Instead, they point to three complementary explanations for the limited political prominence of international redistribution: pluralistic ignorance, low salience, and national political framing.

%LB: Petition - Fabre 25
In a petition experiment, \citet{fabre_majority_2025} find that willingness to sign a petition supporting the GCS is nearly as high as support in the case where respondents are asked directly, and that the gap between the two methods is similar for National Redistribution (NR)\footnote{NR: National Redistribution policy.}. In the United States, 51\% would sign the GCS petition, compared with 53\% who support it in the direct question. In Europe, the corresponding figures are 69\% and 76\%. The differences are comparable for NR: 57\% would sign the petition and 58\% support it directly in the United States, while the corresponding figures are 67\% and 72\% in Europe. Spanish respondents report the highest willingness to sign, at 78\% for the GCS and 74\% for NR \citep[Figure S27]{fabre_majority_2025}.

%LB: List experiment - Fabre 25
\citet{fabre_majority_2025} also conduct a list experiment to assess whether social desirability inflates stated support for the GCS\footnote{\textit{Social desirability bias} occurs when respondents report opinions that present them in a favorable light.}. Respondents report only how many policies they support, while the researchers randomly vary the composition of the list and by including or not the GCS. This design estimates tacit support for the scheme without requiring respondents to endorse it explicitly. The experiment yields an estimated support rate of 62\% overall, including 52\% in the United States and 72\% in Europe. Since these estimates are not significantly different than direct stated support, the authors find no evidence of social desirability bias \citep[Table 1]{fabre_majority_2025_si}. 

%LB: conjoint analysis - Fabre 25 - version étendue 
\citet{fabre_majority_2025} then use a series of conjoint analyses in which respondents choose between policy bundles or candidates' platforms. In the first analysis, respondents choose between a bundle combining the GCS, NR and a national climate policy: (C) (a coal phase-out in the U.S.) and an otherwise identical bundle containing only NR and (C). The bundle including the GCS is preferred by 55\% of U.S. respondents and 74\% of European respondents. 
In a second analysis, respondents are randomly assigned to comparisons in which NR is held constant while the bundles add the GCS, (C), or both. Support for the GCS conditional on NR reaches 55\% in the United States and 77\% in Europe and is not significantly different from support for the GCS alone. Moreover, when respondents choose directly between NR combined with the GCS and NR combined with (C), 53\% of U.S. respondents prefer (C), while 52\% of European respondents prefer the GCS, neither difference is statistically significant. These results indicate that support for the GCS depends neither on compensation through NR nor on its combination with a national climate policy \citep[Figure S8]{fabre_majority_2025}. 
A third analysis presents respondents with progressive and conservative candidates' platforms and randomizes only whether the progressive platform includes the GCS. Adding the GCS increases the probability of being chosen in any country by 2.8 percentage points on average and by 11.2 points in France \citep[Supplementary Table 2]{fabre_majority_2025}.

%===============================================================================================
% - campaign effect (Gustafson et al)
%===============================================================================================

Drawing on two surveys of U.S. voters conducted in December 2018 and April 2019, \citet{gustafson_development_2019} examine how attitudes toward the Green New Deal (GND) changed after the proposal entered the national political debate. Each survey presented respondents with a description of the proposal and asked: ``How much do you support or oppose this idea?'' In December, 81\% supported the GND, including 92\% of Democrats and 64\% of Republicans. By April, the share of voters who had heard at least ``a little'' about the proposal had risen from 17\% to 59\% \citep[Figure 1]{gustafson_development_2019}. Over the same period, support declined among Republicans but remained stable among Democrats. It fell from 64\% to 44\% among Republicans \citep[Figure 2]{gustafson_development_2019}. 
In April, Republican support varied sharply with prior exposure: 85\% of those who had heard nothing about the GND supported it, compared with only 4\% of those who had heard a lot \citep[Figure 3]{gustafson_development_2019}. Among frequent Republican viewers of Fox News, support fell from 54\% to 22\%, compared with a decline from 71\% to 56\% among infrequent viewers. 

%===============================================================================================
% - pluralistic ignorance
%===============================================================================================

Furthermore, \citet{fabre_public_2026} ask respondents whether they accept the Global Climate Scheme (GCS) and then elicit their second-order beliefs about the share of their compatriots who accept it. The median respondent underestimates national acceptance by 16 percentage points, while non-American respondents underestimate acceptance in the United States by 22 points \citep[Figure S30]{fabre_public_2026}. These beliefs correlate strongly with respondents' own attitudes: 51\% of GCS supporters believe that a majority of their compatriots also accept the scheme, compared with only 24\% of opponents. The author reports that ``support for the GCS reaches 72\% among the 39\% of respondents who believe that a majority in their country supports it, compared to 44\% among those who do not'' \citep[Figure S30]{fabre_public_2026}. These associations may reflect both pluralistic ignorance and a ``false consensus effect''.

%===============================================================================================
% - low salience => we should talk about "acceptance" rather than "support"
%===============================================================================================

Before the survey revealed its topic, respondents were randomly assigned one of four broad open-ended questions asking about their main concerns, their needs or wishes, an issue they considered important but neglected in public debate, or the greatest injustice. Among responses to the questions about concerns and wishes, 31\% mentioned ``money'' or the ``cost of living''. Among responses to the greatest-injustice question, only 3.2\% explicitly referred to global inequality, while global redistribution almost never appeared as a wish \citep[Figure 3]{fabre_public_2026}. Taken together, these findings reveal broad acceptance when respondents evaluate a specific proposal but low salience in their spontaneous political concerns. They do not establish that this acceptance translates into active support capable of shaping votes or political mobilization.


%===============================================================================================
%- national framing: original results (https://github.com/bixiou/budget_26/blob/main/figures/group_defended.pdf
%Weighted regression:
%                           Estimate Std. Error t value Pr(>|t|)   
%(Intercept)                  0.16476    0.05966   2.762  0.00586 **
%variant_group_defendedworld  0.16152    0.08268   1.954  0.05103 . 
%---
%Signif. codes:  0 ‘***’ 0.001 ‘**’ 0.01 ‘*’ 0.05 ‘.’ 0.1 ‘ ’ 1
%
%Residual standard error: 1.306 on 998 degrees of freedom
%Multiple R-squared:  0.003809,	Adjusted R-squared:  0.002811 

%Non-weighted:
%                                               Estimate Std. Error t value Pr(>|t|)   
%(Intercept)                                0.15933    0.05952   2.677  0.00755 **
%variant_group_defendedworld  0.20014    0.08230   2.432  0.01520 * 
%---
%Signif. codes:  0 ‘***’ 0.001 ‘**’ 0.01 ‘*’ 0.05 ‘.’ 0.1 ‘ ’ 1

%Residual standard error: 1.3 on 998 degrees of freedom
%Multiple R-squared:  0.005891,	Adjusted R-squared:  0.004895 
%)
%===============================================================================================

A final explanation is a nationalist bias. Fabre in an original survey experiment provides evidence for this mechanism. Respondents were randomly asked either which group they defend when voting or which group they would defend in hypothetical world elections. Responses were coded on a scale ranging from (−2) for ``Family and self'', through \(0\) for ``Fellow citizens'', to (2) for ``Sentient beings''. Framing the elections at the world level shifted responses by 0.162 points toward more universalist groups in the weighted regression, both estimates are significant (\url{https://github.com/bixiou/budget_26/blob/main/figures/group_defended.pdf}). 
